% ===========================================================================
\section{Interfaz de la PC y herramientas}
\label{sec:interfaz}

\subsection{Debugger}

El debugger es una aplicación de texto para la terminal hecha con Textual, en Python
(\repo{tools/debugger/}, decisiones 008 y 040). Se eligió una TUI porque corre igual en Windows y
Linux, no necesita un entorno gráfico y permite mostrar mucha información a la vez. Está dividido
en cuatro capas, de las que solo la última depende de Textual (\repo{docs/spec/interfaz_pc.md}):

\begin{center}
\small
\begin{tabular}{@{}l l@{}}
  \toprule
  \textbf{Capa} & \textbf{Responsabilidad} \\
  \midrule
  \texttt{client}  & Arma y valida tramas, timeouts y reintentos; un método por comando; recibe \texttt{EVT\_HALTED} \\
  \texttt{model}   & Volcado de un instante, latches campo por campo, qué cambió y riesgos inferidos \\
  \texttt{session} & Secuencias de comandos (un paso es \texttt{STEP} + \texttt{READ\_ALL}), historial, controles habilitados \\
  \texttt{tui}     & Paneles, colores y teclas \\
  \bottomrule
\end{tabular}
\end{center}

Así el cliente se prueba con las tramas de ejemplo del protocolo sin terminal, y los scripts de
prueba en placa lo usan sin levantar la interfaz.

\captura[0.45\textheight]{debugger_demo_datos_ciclo08.png}{Debugger en la placa}{Ciclo 8 de
\repo{demo_datos.s}: stall por carga-uso. El pipeline muestra la burbuja en EX y el PC e IF/ID
retenidos; \texttt{sub x4} está en WB escribiendo \texttt{x4} mientras \texttt{add x6} lo lee en
ID.}{fig:debugger}

La pantalla (figura~\ref{fig:debugger}) tiene:
\begin{itemize}[nosep]
  \item \textbf{Programa}: el código desensamblado con la etapa en que está cada instrucción.
  \item \textbf{Pipeline}: las cinco etapas con el PC y la instrucción de cada una, las burbujas
        con su causa y los riesgos del ciclo.
  \item \textbf{Latches}: los cuatro registros de segmentación campo por campo, con los nombres de
        la especificación.
  \item \textbf{Registros} y \textbf{memoria de datos usada}, con lo que cambió en el último paso
        resaltado.
  \item \textbf{Historial}: un renglón por paso con los eventos y los cambios; cualquier paso
        anterior se puede volver a ver.
\end{itemize}

\paragraph{Riesgos inferidos.} La Debug Unit no informa stalls ni forwarding: la PC los deduce de
los campos de los latches con las mismas ecuaciones del hardware y los marca como ``inferidos''
(\decision{043}). Si dos volcados son de ciclos seguidos, además compara lo previsto con lo
observado y marca cualquier diferencia como inconsistencia. Las reglas se prueban contra el modelo
de referencia en cada ciclo de todos los programas de \repo{sw/}.

\begin{center}
\small
\begin{tabular}{@{}l l l@{}}
  \toprule
  \textbf{Tecla} & \textbf{Control} & \textbf{Comandos} \\
  \midrule
  F1          & Conectar               & \texttt{INFO}, \texttt{STATUS}, \texttt{READ\_ALL} \\
  F2 / F8     & Cargar / reprogramar   & ensambla el \texttt{.s} y manda \texttt{LOAD} \\
  F5          & Ejecutar               & \texttt{RUN} y espera \texttt{EVT\_HALTED} \\
  F10         & Paso                   & \texttt{STEP} + \texttt{READ\_ALL} \\
  Shift+F10   & $N$ pasos              & $N$ veces \texttt{STEP} y un \texttt{READ\_ALL} \\
  F4          & Correr hasta el cursor & \texttt{STEP} hasta que la instrucción elegida llega a WB \\
  F6          & Detener                & \texttt{STOP} \\
  F3          & Reiniciar              & \texttt{RESET} \\
  \bottomrule
\end{tabular}
\end{center}

\subsection{Ensamblador}

\repo{tools/assembler} (\texttt{rvasm}) es un ensamblador propio en Python, sin dependencias
(\decision{017}). Acepta el subset de \repo{isa.md}, etiquetas y los nombres ABI de los registros,
y genera \texttt{.hex}, \texttt{.bin}, \texttt{.coe} y un listado con direcciones. Rechaza con
número de línea cualquier instrucción fuera del subset o inmediato fuera de rango. Lo usan el
debugger, el modelo de referencia y los scripts de verificación. Sus salidas se comparan palabra por
palabra con las de un ensamblador estándar (GNU o LLVM, \repo{validate_reference.py}).

\subsection{Modelo de referencia}

\repo{tools/isasim} es un simulador en Python escrito aparte del RTL (\decision{020}), en dos capas:

\begin{itemize}[nosep]
  \item \textbf{ISA}: el efecto de cada instrucción, decodificada por su cuenta desde
        \repo{isa.md}, sin usar las tablas del ensamblador.
  \item \textbf{Pipeline}: en qué etapa está cada instrucción en cada ciclo, con las reglas de
        \repo{pipeline.md} (forwarding, stall, flush y parada por HALT), y qué tiene cada
        registro de segmentación.
\end{itemize}

Produce el estado final, la cantidad de ciclos y la traza de los cuatro latches empaquetados en
cada ciclo. Esa traza es la que compara el banco de pruebas del RTL y también la que compara
\repo{comparar.py} con la placa (sección~\ref{sec:verif}).
